On veut préparer $\qty{100}{\mL}$ de solution de chlorure de fer(III) de
concentration molaire $\qty{0.250}{\mol\per\L}$.
\begin{enumerate}
    \item Quelle est la formule de chlorure de fer(III)~?
    \item Calculer la quantité de matière de soluté que l'on doit dissoudre pour
          préparer cette solution~; en déduire la masse de chlorure de fer(III)
          qu'il faut peser.
    \item Écrire l'équation de la réaction de dissolution du chlorure de
          fer(III). En déduire les concentrations des ions présents dans
          celle-ci.
\end{enumerate}
\begin{donnees}
    $M_{\ce{Fe}}=\qty{55.8}{\g\per\mol}$~;
    $M_{\ce{Cl}}=\qty{35.5}{\g\per\mol}$.
\end{donnees}

